% Requires \usepackage{graphicx,booktabs}. Paths assume this figures/ directory is beside the main TeX file.

\begin{table*}[t]
\centering
\small
\caption{Controlled Task1 continuation from historical SGD+EWC and visual GAM+EWC Task0 anchors, with SGD+EWC used for both continuations (training seed 0). The class-incremental accuracies $a_{11},a_{21},a_{22}$ (\%) use all classes seen at the evaluation stage: old-task accuracy before and after continuation, and new-task accuracy after continuation, respectively. $a_{11}$ is the current FP32 anchor evaluation; $a_{21}$ and $a_{22}$ are from the continuation accuracy matrix. $D_0$ is signed old-task cross-entropy change relative to each anchor, measured separately with the fixed Task0 20-class head and deterministic preprocessing on 1,259 train and 307 test examples. Negative damage means loss decreased. Raw trace at Task0 averages squared per-example unscaled CE gradients over the same 256 old-task training examples in the ambient coordinates of 24 effective Q/V $BA$ matrices (scaling 1), without normalization, floor, or regularization. Displacement is the Euclidean/Frobenius norm in those effective-weight coordinates. The inherited EWC Fisher matrices may differ across anchors; this is not an isolated protection-matrix intervention. The new-task plasticity guard is retained: GAM-anchor $a_{22}$ differs by $-2.75$ percentage points from the SGD-anchor row.}
\label{tab:rq2_controlled_continuation}
\begin{tabular}{lrrrrrrr}
\toprule
Task0 anchor & $a_{11}$ & $a_{21}$ & $a_{22}$ & $D_0^{\rm train}$ & $D_0^{\rm test}$ & $\mathrm{tr}(\widehat F_0)$ & $\|\Delta_1\|_2$ \\
\midrule
SGD+EWC & 85.6678 & 84.36 & 74.80 & +0.003674 & -0.019246 & 1465.497 & 0.374950 \\
GAM+EWC & 87.6221 & 84.04 & 72.05 & -0.002157 & -0.013388 & 761.901 & 0.349876 \\
\bottomrule
\end{tabular}
\end{table*}


\begin{figure*}[t]
\centering
\includegraphics[width=0.95\textwidth]{figures/Figure_B_neighborhood.pdf}
\caption{Equal-radius neighborhoods around historical SGD+EWC and visual GAM+EWC Task0 anchors (training seed 0; current full-FP32 deterministic re-evaluation). For each split and anchor, define $D(r,u)=L_0(\theta_0+ru)-L_0(\theta_0)$ and $D_+=\max(D,0)$. (a) $G_{\rm mean}=K^{-1}\sum_k D_+(r,u_k)$. (b) $G_{95}$ is the empirical 95th percentile of $D_+$ (linear interpolation). Solid and dashed curves denote train and test, respectively; train/test losses use 1,259/307 old-task examples and the fixed Task0 20-class head. (c) Mean, empirical 95th percentile, and sampled maximum of the raw empirical Fisher trace across directions, where $\widehat F_{\rm emp}=N^{-1}\sum_{i=1}^{N}g_i g_i^\top$, $N=256$, and $g_i$ is the gradient of the unscaled per-example CE on the same stratified old-task training examples. The empirical mean includes $1/N$; raw means no trace normalization, floor, or regularization. All perturbations and gradients use the ambient coordinates of 24 effective Q/V $BA$ matrices (scaling 1); $r$ is their joint Euclidean/Frobenius displacement. Each positive radius has $K=16$ paired unit directions shared across anchors; $r=0$ is one reused origin measurement. These are direction statistics, not independent training replicates, confidence intervals, or worst-case neighborhood guarantees. Radius axes in (a--c) use a symmetric-log scale with a linear region through $r=0.5$; damage axes in (a,b) use a symmetric-log scale with a linear region through $10^{-4}$. (d) Prefix-safe $G_{\rm mean}$ tolerance is the largest sampled radius such that all sampled radii up to it meet threshold $\epsilon$. Co-located blue outer and orange inner markers/lines preserve exact ties without moving data; triangles and $\geq20$ indicate right censoring at the largest tested radius. Both methods tie at every tested $G_{\rm mean}$ tolerance threshold on each split. GAM+EWC has smaller train $G_{\rm mean}$ and mean raw trace on the measured grid, while its test $G_{\rm mean}$ is larger at every positive radius; the plot does not support a higher $G_{\rm mean}$ tolerance radius.}
\label{fig:rq2_neighborhood}
\end{figure*}

\begin{figure*}[t]
\centering
\includegraphics[width=0.95\textwidth]{figures/Figure_C_sequence_and_endpoint.pdf}
\caption{Historical SGD+EWC and visual GAM+EWC sequences, training seed 0, re-evaluated in current full FP32. The historical sequences use SGD-to-SGD and GAM-to-GAM optimizers, respectively, so their future optimizers differ; Table A separately controls the future optimizer, including GAM-to-SGD. (a) Task0 test accuracy and (b) signed Task0 test CE damage relative to each method's Task0 anchor after completed tasks $t=0,\ldots,9$, using the fixed local 20-class head and the same 307 deterministic test examples throughout; these are not all-seen-class accuracies. (c) Measured endpoint directional sensitivity only at $t=9$, $S_{\rm path}(9)=I_9/\|\Delta_9\|_2^2$, where $I_9=\int_0^1\Delta_9^\top\widehat F(\theta_0+s\Delta_9)\Delta_9\,ds$. The empirical Fisher is the mean of per-example outer products of unscaled CE gradients over the same 256 stratified old-task training examples, with no trace normalization, floor, or regularization and with the Task0 head fixed. The effective-weight straight segment between the Task0 and Task9 endpoints is evaluated at 9 composite-Simpson nodes; it is not the actual optimizer-step training trajectory. Both 5-versus-9-node quadrature checks pass the 5\% gate and the matched-subset numerical bound checks pass the stated 2\% tolerance. No $S_{\rm path}$ measurements are available at $t=1,\ldots,8$, so no sensitivity trajectory is drawn. (d) Displacement $\|\Delta_t\|_2$ is measured in the ambient coordinates of 24 effective Q/V $BA$ matrices (scaling 1). GAM+EWC has lower endpoint sensitivity despite a larger endpoint displacement. Curves connect task checkpoints as visual guides; this single-seed case supplies no uncertainty estimate over training runs.}
\label{fig:rq2_sequence_endpoint}
\end{figure*}

% Appendix tables (move to the appendix when integrating):

\begin{table}[t]
\centering
\small
\caption{Appendix: prefix-safe tolerance radii based on the empirical 95th percentile of positive Task0 CE damage, $G_{95}$, for historical SGD+EWC and visual GAM+EWC anchors. Each positive radius uses the same 16 paired directions; these are direction statistics from training seed 0, not training replicates. For each threshold, the reported radius is the largest sampled radius such that every sampled radius up to it passes the threshold. The sampled grid is $r\in\{0,0.5,1,2,5,10,20\}$. $\geq 20$ denotes right censoring at the grid boundary, not a measured radius above 20. Both deterministic train and test losses use the fixed Task0 20-class head. These are finite-direction, finite-grid measurements, not worst-case neighborhood guarantees.}
\label{tab:rq2_q95_tolerance}
\begin{tabular}{lrrr}
\toprule
Split & $\epsilon$ & SGD+EWC & GAM+EWC \\
\midrule
Test & 0.001 & 0 & 0 \\
Test & 0.005 & 1 & 1 \\
Test & 0.01 & 2 & 2 \\
Test & 0.02 & 5 & 5 \\
Test & 0.05 & 10 & 10 \\
Train & 0.001 & 1 & 2 \\
Train & 0.005 & 5 & 5 \\
Train & 0.01 & 5 & 5 \\
Train & 0.02 & 10 & 10 \\
Train & 0.05 & 10 & $\geq 20$ \\
\bottomrule
\end{tabular}
\end{table}


\begin{table*}[t]
\centering
\small
\caption{Appendix: existing two-anchor by two-realized-direction bridge (training seed 0). SGD and GAM abbreviate historical SGD+EWC and visual GAM+EWC checkpoints. Each direction is the unit effective-weight displacement from Task0 to Task9 of its named historical sequence, transplanted onto the named Task0 anchor; $r$ is the absolute Euclidean radius in ambient coordinates of the 24 effective Q/V $BA$ matrices (scaling 1). The signed train/test Task0 CE changes use each anchor's fixed 20-class head and deterministic preprocessing. The raw, unnormalized empirical Fisher trace uses per-example unscaled CE gradients on the shared 256 training examples. Radius zero reuses the anchor measurement for both directions. This preserves the existing historical-sequence direction identity: these directions are not the new controlled-continuation directions of Table~\ref{tab:rq2_controlled_continuation}.}
\label{tab:rq2_existing_bridge}
\begin{tabular}{llrrrr}
\toprule
Anchor & Direction & $r$ & $D^{\rm train}$ & $D^{\rm test}$ & $\mathrm{tr}(\widehat F)$ \\
\midrule
SGD & SGD & 0 & +0.000000 & +0.000000 & 1465.497 \\
SGD & SGD & 0.25 & -0.000398 & -0.000103 & 1554.181 \\
SGD & SGD & 0.5 & +0.010850 & +0.016222 & 1799.319 \\
SGD & SGD & 1 & +0.085561 & +0.147516 & 3209.145 \\
SGD & SGD & 1.5 & +0.332180 & +0.517852 & 7572.045 \\
SGD & SGD & 2 & +0.970554 & +1.276820 & 15993.630 \\
SGD & GAM & 0 & +0.000000 & +0.000000 & 1465.497 \\
SGD & GAM & 0.25 & +0.002023 & -0.001218 & 1595.781 \\
SGD & GAM & 0.5 & +0.006002 & +0.001458 & 1600.829 \\
SGD & GAM & 1 & +0.020459 & +0.018917 & 1779.211 \\
SGD & GAM & 1.5 & +0.044087 & +0.051845 & 1989.137 \\
SGD & GAM & 2 & +0.080304 & +0.100854 & 2162.101 \\
GAM & SGD & 0 & +0.000000 & +0.000000 & 761.901 \\
GAM & SGD & 0.25 & +0.000266 & -0.008249 & 806.176 \\
GAM & SGD & 0.5 & +0.006800 & +0.000914 & 985.959 \\
GAM & SGD & 1 & +0.056009 & +0.111834 & 2067.648 \\
GAM & SGD & 1.5 & +0.232725 & +0.455334 & 5270.921 \\
GAM & SGD & 2 & +0.698221 & +1.115471 & 13016.053 \\
GAM & GAM & 0 & +0.000000 & +0.000000 & 761.901 \\
GAM & GAM & 0.25 & +0.002057 & -0.002281 & 724.649 \\
GAM & GAM & 0.5 & +0.005642 & -0.000945 & 767.457 \\
GAM & GAM & 1 & +0.017987 & +0.013447 & 950.242 \\
GAM & GAM & 1.5 & +0.039360 & +0.041232 & 1175.963 \\
GAM & GAM & 2 & +0.071637 & +0.088605 & 1340.090 \\
\bottomrule
\end{tabular}
\end{table*}

